\documentclass[a4paper,9pt]{extarticle}

%=========================================================
% PACKAGES
%=========================================================

\usepackage[
a4paper,
top=1.8cm,
bottom=1.5cm,
footskip=7mm,
left=1cm,
right=1cm,
headheight=14pt,
headsep=5mm
]{geometry}
\usepackage[T1]{fontenc}
\usepackage[utf8]{inputenc}
\usepackage[french]{babel}

\usepackage[sfdefault]{FiraSans}
\renewcommand{\familydefault}{\sfdefault}

\usepackage{multicol}
\setlength{\columnsep}{0.7cm}

\usepackage{amsmath,amssymb}

\usepackage{titlesec}
\usepackage{enumitem}

\usepackage{tcolorbox}
\tcbuselibrary{skins,breakable}

\usepackage{xcolor}

\usepackage{graphicx}
\usepackage{float}

\usepackage{tikz}

\usepackage{fancyhdr}
\usepackage[hidelinks]{hyperref}


%=========================================================
% PAGE SETTINGS
%=========================================================

\setlength{\parindent}{0pt}
\setlength{\parskip}{1.5pt}

\setlength{\abovedisplayskip}{4pt}
\setlength{\belowdisplayskip}{4pt}

\setlist[itemize]{
leftmargin=*,
itemsep=1pt,
topsep=1pt
}


%=========================================================
% COLORS
%=========================================================

\definecolor{Blue}{HTML}{E3F0FF}
\definecolor{Green}{HTML}{E3F7EA}
\definecolor{Pink}{HTML}{FFE5ED}
\definecolor{Orange}{HTML}{FFF0DC}
\definecolor{Purple}{HTML}{EEE7FF}
\definecolor{Yellow}{HTML}{FFFADB}


\definecolor{BlueTitle}{HTML}{285B8F}
\definecolor{GreenTitle}{HTML}{287052}
\definecolor{PinkTitle}{HTML}{A84063}
\definecolor{OrangeTitle}{HTML}{A96500}
\definecolor{PurpleTitle}{HTML}{6245A5}


%=========================================================
% HEADER
%=========================================================

\pagestyle{fancy}

\fancyhf{}

\fancyhead[L]{\small\textbf{Numerical Methods}}
\fancyhead[C]{\small Admin and Floating Point}
\fancyhead[R]{\small Week 01 --- 07-13/09/2026}

\fancyfoot[C]{\small\href{https://github.com/AlphaXT38G/notes_ba5}{github.com/AlphaXT38G/notes\_ba5}}
\fancyfoot[R]{\small \thepage}

\renewcommand{\headrulewidth}{0.4pt}



%=========================================================
% TITLES
%=========================================================

\titleformat{\section}
{\large\bfseries\color{BlueTitle}}
{}{0pt}{}

\titlespacing{\section}
{0pt}{6pt}{3pt}


% separator

\newcommand{\separator}{
\vspace{2pt}
\hrule
\vspace{3pt}
}


%=========================================================
% BOXES
%=========================================================

% Boxes are kept whole (unbreakable): inside multicols, a box split by tcolorbox
% leaves an empty column behind it.

% DEFINITION - blue
\newtcolorbox{definition}{
enhanced,
unbreakable,
colback=Blue,
colframe=BlueTitle,
boxrule=0.7pt,
arc=0mm,
title=\textbf{Definition},
fonttitle=\bfseries
}

% IMPORTANT - pink
\newtcolorbox{important}{
enhanced,
unbreakable,
colback=Pink,
colframe=PinkTitle,
boxrule=0.7pt,
arc=0mm,
title=\textbf{Key idea},
fonttitle=\bfseries
}

% FORMULAS - green
\newtcolorbox{formula}{
enhanced,
unbreakable,
colback=Green,
colframe=GreenTitle,
boxrule=0.7pt,
arc=0mm,
title=\textbf{Key formulas},
fonttitle=\bfseries
}

% WARNING - orange
\newtcolorbox{warning}{
enhanced,
unbreakable,
colback=Orange,
colframe=OrangeTitle,
boxrule=0.7pt,
arc=0mm,
title=\textbf{Warning},
fonttitle=\bfseries
}

% EXAMPLE - purple
\newtcolorbox{example}{
enhanced,
unbreakable,
colback=Purple,
colframe=PurpleTitle,
boxrule=0.7pt,
arc=0mm,
title=\textbf{Example},
fonttitle=\bfseries
}

% SUMMARY - orange
\newtcolorbox{summary}{
enhanced,
unbreakable,
colback=Yellow,
colframe=OrangeTitle,
boxrule=0.7pt,
arc=0mm,
title=\textbf{Summary},
fonttitle=\bfseries
}



%=========================================================
% FIGURE BOX
%=========================================================

\newtcolorbox{figurebox}{
enhanced,
colback=white,
colframe=gray,
boxrule=0.5pt,
arc=0mm,
center
}



%=========================================================
% DOCUMENT
%=========================================================

\begin{document}

%\footnotesize   --> To reduce the size of the text

\thispagestyle{fancy}



% Columns keep their natural height: boxes are unbreakable, stretching would open gaps.
\raggedcolumns
\begin{multicols}{2}


%=========================================================
\section{Course Organisation}


\begin{summary}

\begin{itemize}

\item \textbf{CS-328}, Prof.\ Wenzel Jakob (Realistic Graphics Lab, RGL).

\item \textbf{Lecture}: Thursday, CE13, 10\string:15-12\string:00 (concepts, live examples,
general Q\&A).

\item \textbf{Exercise session}: Thursday, BC07/08, 16\string:15-18\string:00 (first 3 weeks:
tutorials, then homework Q\&A; ignore the INF3 room).

\item TAs: Dongsu Zhang, Ekrem Yilmazer (PhD); Jordan Warne, Felix Burchardt, Arthur Herbette
(AEs).

\end{itemize}

\end{summary}


\begin{formula}

\textbf{Grading}

\begin{itemize}

\item 4 homeworks during the semester: \textbf{10\%} of the final grade.

\item Pen \& paper exam during the January exam session: \textbf{90\%} of the final grade.

\item Optional practice midterm: graded but 0\% weight.

\item Homework and exam graded on a 1-6 scale, then averaged with these weights (mapping from
raw percentages generally uses a curve).

\end{itemize}

\end{formula}


\textbf{Homeworks}

\begin{enumerate}

\item Basics: NumPy, SciPy, Matplotlib, floating point arithmetic, error analysis, first steps
with a Monte Carlo method.

\item Processing a RAW image: demosaic \& process RAW Bayer grid data from a Canon DSLR, solve a
least squares problem from ``Native Camera'' to human (sRGB) color space.

\item Inverse kinematics: simple IK ``chain'' (a worm body being pulled), SVD to solve a
rank-deficient linear system.

\item Autodiff \& machine learning: build a tiny ML framework (``MyTorch'') and train a tiny
neural network.

\end{enumerate}

Plus 6-7 ungraded exercise sheets. No text book (optional readings); recommended: \emph{Essence
of Linear Algebra} (3Blue1Brown).


\begin{warning}

\begin{itemize}

\item Late policy: 25\% penalty for each extra started late day (no points after 4 days).

\item Survey course: many methods, none in depth. Not a programming language course.

\item Any modern LLM can solve all problems of this course, but reviewing LLM output does not
work to learn: ``to learn, you must struggle''. Huge drop in exam performance in 2026; the exam
will not be curved up this year.

\end{itemize}

\end{warning}


\textbf{Python}: reference version \textbf{3.14} (standard Python or Conda); an interpreted
language with an unusual indentation-based syntax. ``Jupyter notebook'' means \textbf{Jupyter
Lab}.



%=========================================================
\section{What is a Numerical Method?}


\begin{important}

Computer graphics, computer vision, machine learning, computational photography, weather
simulation, geometric information systems\dots: \textbf{all of them have some kind of numerical
method at their core}, running on number crunchers (CPUs, GPUs).

\end{important}


\begin{example}

\begin{itemize}

\item \textbf{Rendering}: visual effects $\simeq$ solving integrals,
$\text{pixel color} = \int \text{light}\,\mathrm{d}x$.

\item \textbf{Fluid simulation}: to compute the next frame (``timestep''), solve an enormous
linear system $\mathbf{x} = \mathbf{A}^{-1}\mathbf{b}$ (HUGE matrix).

\item \textbf{Character animation}: attaching a virtual body to a skeleton using optimization,
$\mathbf{x} = \mathbf{A}^{+}\mathbf{b}$ (pseudoinverse).

\item \textbf{Computer vision}: quickly process images using learned convolutions (``CNNs''); a
convolution is a particular type of linear transformation.

\item \textbf{Machine learning}: computing derivatives of software,
$\frac{\partial}{\partial \mathbf{x}} f(\mathbf{x})$ where $f$ is a computer program.

\end{itemize}

\end{example}


\begin{definition}

\textbf{CS-328}

About the \textbf{computer implementation} of abstract concepts from MATH-111 (Linear Algebra:
$\mathbf{x} = \mathbf{A}^{-1}\mathbf{b}$) and MATH-101/106 (Analysis:
$\frac{\partial}{\partial \mathbf{x}} f(\mathbf{x})$, $y = \int_\Omega f(x)\,\mathrm{d}x$).

Typical computer science works with $\mathbb{N}$ (graph theory, databases, network protocols,
cryptography); numerical methods work with $\mathbb{R}^n$.

\end{definition}


\begin{summary}

\textbf{Many layers} (top to bottom)

\begin{itemize}

\item Applications

\item Frameworks: NumPy, PyTorch, TensorFlow, JAX

\item Optimization algorithms: gradient descent, Newton's method in N dimensions, Broyden's
method, Gauss-Newton, Levenberg-Marquardt

\item Matrix factorizations: the ``Big 3'' LU, QR, SVD

\item Compilation: of floating point math, just-in-time (JAX, XLA, etc.)

\item Floating point arithmetic: history, IEEE-754, error propagation (\emph{today})

\item Processor architecture: superscalar processors, vectorization, memory hierarchy, shared
memory (GPUs)

\end{itemize}

\end{summary}



\end{multicols}


\end{document}
